\section{Method and the pre-registered test}\label{sec:method}

\subsection{The stretch}

Each contest enters once, seen from the favourite's side. Write $p$ for the favourite's scaled price at the snapshot,
between 0.5 and 0.95, and let the favourite either win or lose. The model gives the favourite the chance
\begin{equation}\label{eq:stretch}
  \frac{p^{\,b}}{p^{\,b} + (1-p)^{b}},
\end{equation}
where $b$ is a single number, the \emph{stretch}, shared by every contest in a group. When $b = 1$ the chance equals the
price and the market is calibrated. When $b$ is above 1, the model pushes every price away from 50\%: favourites are
given a higher chance than their price, and the market is \emph{underconfident}. When $b$ is below 1, prices are
pulled towards 50\% and the market is \emph{overconfident}. To give a sense of scale, a stretch of 1.65 turns an 80\%
favourite into a 91\% chance, and a stretch of 0.32 turns it into a 61\% chance.

Equivalently, the stretch multiplies the log-odds of the price: the model's log-odds of a favourite winning are $b$
times $\log\{p/(1-p)\}$. It is therefore a logistic regression of the result on the log-odds of the price with no
intercept, the recalibration slope of \citet{le2026} without its intercept. There is no intercept because the two sides
of a contest are mirror images: the model treats a price of 70\% for one candidate and 30\% for the other the same way
whichever candidate is listed first, so there is no direction in which an intercept could push. We estimate $b$ by
maximum likelihood, the value that gives the actual winners the highest joint chance, using Newton's method. A contest
priced at exactly 50\% says nothing about $b$. When every favourite in a group won, the likelihood keeps rising as $b$
grows and the estimate has no upper limit; we report this as ``no limit'' rather than as a number.

The stretch is deliberately simple. It summarises a whole group of contests with one number that has a direct reading,
it is the correction a forecaster would apply, and it has one parameter to estimate from samples of a few dozen to a
few hundred contests. It assumes that the bias has the same form at every price level, an assumption the calibration
charts in Section~\ref{sec:results} allow the reader to check.

\subsection{Intervals}

Contests are not independent. On a general-election night, a national swing moves races in every state the same way;
within a state, races share voters, polls and polling errors. Treating each contest as an independent draw would make
the estimates look more certain than they are. Every interval in this paper therefore comes from resampling whole
groups of contests: 2{,}000 resamples, drawn with a fixed seed (20261101), each refitting the stretch on the resampled
contests. The 95\% interval runs from the 2.5th to the 97.5th percentile of the 2{,}000 refitted values. When more than
2.5\% of resamples have every favourite winning, the upper end has no limit.

Two resampling units are used. For the training data, which span many election days, the main interval resamples
whole \emph{election days}, so that a night with a common polling miss enters or leaves as one block. For a single
night such as the midterms, days cannot be resampled, and the interval resamples whole \emph{states}. We also report
state-resampled intervals for the training data, so that every group can be compared on the same footing. Resampling
single contests, which ignores both kinds of dependence, is reported only as a reference.

\subsection{The pre-registered test}

The test was written, run on the 2024 data as a check, and archived before the vote. It applies to each platform
separately.

\paragraph{Sample.} US general-election contests voted on 3 November 2026 that are on the list made before the vote,
have settled with one of their two listed candidates winning, and pass every price rule, priced at 22:00 UTC on
3 November. Every office counts. Contests in Georgia, Louisiana and Mississippi are excluded, because a race there can
be decided by a later runoff rather than by the 3 November vote. A platform with fewer than five contests is not
tested.

\paragraph{Two questions.} The test fits the stretch on the midterm sample and computes its 95\% interval by resampling
states. It then asks:
\begin{itemize}
  \item \textbf{Level.} If the whole interval lies above 1, the platform's general-election markets are
        \emph{underconfident}; if it lies below 1, \emph{overconfident}. Otherwise there is \emph{no detectable
        miscalibration}.
  \item \textbf{Change.} The same interval is compared with the platform's stretch on 5 November 2024, fixed in
        advance at 1.972 for Kalshi and 1.828 for Polymarket. If the whole interval lies below that value, the
        stretch is \emph{lower than 2024}; if above it, \emph{higher than 2024}. Otherwise it is \emph{consistent with
        2024}.
\end{itemize}
The 2024 value is treated as a fixed number, not as an estimate with its own interval. This was decided before the vote,
on the basis of a check on the 2024 data: comparing two intervals of this width cannot detect even a large change, and
the question of interest is whether the 2026 markets differ from what was seen in 2024. A stretch with no upper limit is
kept as it is, and the verdicts read only the ends of the interval.

\paragraph{Readings.} The two answers combine into one of four readings, fixed in advance (Table~\ref{tab:readings}).
The first supports the explanation that primaries are traded differently from general elections; the second and third
support, fully or in part, the explanation that the underconfidence of 2024 faded over time.

\begin{table}[t]
\centering\small
\caption{The pre-registered readings of the midterm test, per platform.}\label{tab:readings}
\begin{tabular}{p{4.6cm}p{3.3cm}p{6.2cm}}
\toprule
Level & Change & Reading \\
\midrule
Underconfident & Consistent with or higher than 2024 & Underconfidence persists in general elections: supports ``primaries are traded differently''. \\
No detectable miscalibration, or overconfident & Lower than 2024 & Underconfidence has not persisted: supports ``it switched off over time''. \\
Underconfident & Lower than 2024 & Still underconfident, but less than in 2024: partial support for ``it switched off over time''. \\
Any other combination & & Inconclusive. \\
\bottomrule
\end{tabular}
\end{table}

\paragraph{Stages.} Midterm markets settle over days or weeks, and some only once results are certified. The test
therefore runs twice. A \emph{preliminary} run, allowed from 10 November 2026, uses the contests settled by then and
reports how many listed contests are still open. Races still uncalled are more often those in which the favourite
underperformed, so the preliminary run leans towards finding underconfidence. The \emph{final} run takes place once
every listed contest has settled, or on 31 January 2027, whichever comes first; contests still unsettled then are left
out. The reading is judged on the final run only. This version of the paper reports the preliminary run.

\paragraph{Safeguards.} Before testing, the script recomputes the 2024 stretch for each platform from the price table
and stops if it differs from the frozen value. That guarantees the midterm rows were added without changing any earlier
data. It also stops if any midterm contest is priced at a moment other than the fixed snapshot.

\subsection{What was frozen, and where}

Before 3 November 2026, the complete analysis was made public and archived:
\begin{itemize}
  \item the code, the training results and the midterm test, as release \texttt{v1.0-preregistration} of the public
        repository, archived on Zenodo \citep{kirillov2026code};
  \item the raw market data downloaded from both exchanges, archived separately \citep{kirillov2026data};
  \item a list of SHA-256 fingerprints of the 131 data and result files, so that anyone can check that a later copy
        matches the archived one byte for byte;
  \item the list of midterm contests, made from the open markets before the vote.
\end{itemize}
The test script was run on 5 November 2024 as a self-test and reproduced the frozen 2024 values. Code written after the vote to
download and clean the settled midterm markets, and one change to where Polymarket trades are fetched from, are
described in Section~\ref{sec:limitations}. None of them changes a rule.

\subsection{Analyses beyond the pre-registered test}

The following analyses were not part of the frozen test. They are reported as exploratory, and where they produce
p-values, these are descriptive.

\paragraph{Forecast accuracy.} Prices and stretched prices are scored by their \emph{log loss}, the average of minus
the log of the chance given to what actually happened, and by their \emph{Brier score}, the average squared gap
between the chance and the result; lower is better on both. In the training data, each election day is scored with a
stretch fitted on all the other days. On the midterms, the stretch values are fixed before the vote: the frozen 2024
value, and the values fitted on all earlier general elections, all earlier primaries and all earlier contests. The
improvement is the log loss of the raw prices minus that of the stretched prices, with an interval from resampling
states while the stretch is held fixed. As a reference, the midterms are also scored with their own fitted stretch,
which is the best any single stretch could have done.

\paragraph{Significance.} For the level question, we report the share of state resamples at or below 1, and a test of
the stretch against 1 whose standard error allows for contests in one state moving together. For the accuracy
question, we compare the log loss of raw and stretched prices with a standard error that allows for the same
dependence. Naive versions that treat contests as independent are shown alongside, to make the effect of this
dependence visible.

\paragraph{Same races on both platforms.} The list made before the vote names each race (for example, the House seat
for a given district). Races that both platforms listed and that passed every rule on both are matched by that name,
and the stretch is fitted again on this shared set on each platform. This separates differences in how the platforms
priced the same races from differences in which races they listed.

\paragraph{Other checks.} The training results are recomputed for all countries instead of the US only, and on Kalshi
without the five contests in which neither candidate traded in the day before the snapshot (Polymarket's rules
already exclude such contests; Kalshi's do not). Primaries are split into states with and without runoffs, and by
how long their markets took to settle. The trading backtest is described in Section~\ref{sec:backtest}.
